% Chern-Simons Implementation Notes
% Source: docs/chern-simons-notes.md (migrated to LaTeX)
% Uses macros from preamble.tex

\section{Chern-Simons Implementation Notes}
\label{sec:chern-simons}

\subsection{About This File}

This is an \textbf{example-specific} implementation guide for the Chern-Simons 2+1D case.

\paragraph{When to create similar files.}
\label{par:cs-when-to-create}
\begin{itemize}
  \item For any complex example requiring hybrid approaches (symbolic + manual)
  \item When implementing topological field theories with special tensor structures
  \item For examples with gauge symmetries or constraint equations
  \item When exploring new physics domains (Yang-Mills, BF theory, supergravity, etc.)
\end{itemize}

\paragraph{Template for new example notes.}
\label{par:cs-template}
\begin{enumerate}
  \item Physics Background (Lagrangian, EOM, physical interpretation)
  \item Implementation Status (what works, what doesn't, current limitations)
  \item Wolfram Implementation Pattern (specific approaches for this example)
  \item JSON Structure (how the equations are represented)
  \item Python Simulation (grid setup, initial conditions, expected physics)
  \item Future Automation Plan (steps to fully automate if currently hybrid)
  \item References (scripts, JSON files, related examples)
\end{enumerate}

\paragraph{Integration with main docs.}
\label{par:cs-integration}
\begin{itemize}
  \item Reference this file from \texttt{MEMORY.md} under ``Example Implementations''
  \item Add error patterns to \texttt{troubleshooting.md} if they're general enough
  \item Link back to main docs for general patterns (this file is for specifics)
\end{itemize}

\subsection{Physics Background}
\label{sec:cs-physics}

Chern-Simons theory in 2+1D:
\begin{equation}
  \lag = -\tfrac{1}{4}\, F_{\mu\nu}\, F^{\mu\nu}
       + \frac{\kappa}{2}\, \varepsilon^{\mu\nu\rho}\, A_\mu\, \partial_\nu A_\rho
  \label{eq:cs-lagrangian}
\end{equation}

Equations of motion:
\begin{equation}
  \partial_\mu F^{\mu\nu} + \kappa\, \varepsilon^{\nu\alpha\beta}\, \partial_\alpha A_\beta = 0
  \label{eq:cs-eom}
\end{equation}

In Lorenz gauge ($\partial^\mu A_\mu = 0$):
\begin{align}
  \ddt^2 A_0 &= \lapl A_0 + \kappa\bigl(\ddx A_2 - \ddy A_1\bigr) \label{eq:cs-lorenz-A0} \\
  \ddt^2 A_1 &= \lapl A_1 + \kappa\,\ddy A_0 \quad\text{[time term $-\kappa\,\ddt A_2$ moves to LHS]} \label{eq:cs-lorenz-A1} \\
  \ddt^2 A_2 &= \lapl A_2 - \kappa\,\ddx A_0 \quad\text{[time term $+\kappa\,\ddt A_1$ moves to LHS]} \label{eq:cs-lorenz-A2}
\end{align}

\subsection{Implementation Status}
\label{sec:cs-status}

\subsubsection{What Works (Fully Automated)}
\label{sec:cs-works}

\begin{itemize}
  \item 3D manifold setup (M3, 2+1D spacetime)
  \item Full Maxwell-Chern-Simons Lagrangian symbolic derivation
  \item \textbf{Automatic epsilon tensor evaluation} (\texttt{EvaluateEpsilonComponents})
  \item Component decomposition for vector fields in 3D
  \item \textbf{Automatic gradient direction detection} (\texttt{gradient\_x} vs \texttt{gradient\_y})
  \item \textbf{Cross-field reference detection} ($A_0 \to A_1 \to A_2$)
  \item Python simulation with gradient operators
  \item Energy transfer between components (CS coupling effect)
\end{itemize}

\subsubsection{Known Limitations}
\label{sec:cs-limitations}

\begin{itemize}
  \item Mixed time-space derivatives ($\ddt\,\ddx$) classified as ``laplacian'' in JSON
  \item Extra cross-Laplacian terms appear (from unsimplified gauge terms)
  \item Time derivative separation pattern needs 2+1D update for cleaner JSON
\end{itemize}

\subsection{Wolfram Implementation Pattern}
\label{sec:cs-wolfram}

\subsubsection{Fully Automated Approach (Current)}
\label{sec:cs-automated}

\begin{lstlisting}[style=mathematica]
(* 1. Build full Lagrangian with epsilon tensor *)
MaxwellLagrangian = -1/2 CD3[-a][csA[-b]] eta3[a,c] eta3[b,d] CD3[-c][csA[-d]]
CSLagrangian = (kappa/2) * epsiloneta3[a, b, c] * csA[-a] * CD3[-b][csA[-c]]
FullLagrangian = MaxwellLagrangian + CSLagrangian

(* 2. Derive EOM - includes both Maxwell and CS contributions *)
eomFull = VarD[csA[-a], CD3][FullLagrangian]

(* 3. Decompose - EvaluateEpsilonComponents handles epsilon automatically *)
componentEqs = DecomposeToComponents[eomFull, csA[-a], cart3]

(* 4. Export - IdentifyGradientDirection detects x/y correctly *)
jsonStructure = BuildMultiFieldJSONStructure[fullEquations, metadata]
\end{lstlisting}

\subsubsection{Previous Hybrid Approach (No Longer Needed)}
\label{sec:cs-hybrid}

The old approach required manual CS term addition after decomposition.
This is no longer necessary with the automated epsilon evaluation.

\subsubsection{Why Automation Now Works}
\label{sec:cs-why-automation}

The epsilon tensor $\varepsilon^{\mu\nu\rho}$ creates complex index structures:

\begin{itemize}
  \item xAct's \texttt{epsiloneta3[-a,-b,-c]} exists but component evaluation is non-trivial
  \item \texttt{ToBasis[chart][epsilon[...]]} produces symbolic basis components
  \item Need explicit rules for all index permutations with metric signature factors
  \item Simpler to compute CS terms from known structure than fully automate
\end{itemize}

\subsubsection{Attempted (Failed) Approaches}
\label{sec:cs-failed}

\paragraph{Direct epsilon in Lagrangian.}
\label{par:cs-direct-epsilon}
\begin{lstlisting}[style=mathematica]
(* Symbolic epsilon remains after decomposition *)
CSTerm = epsilon[a,b,c] A[-a] CD[-b][A[-c]]
(* Result: eps[{0,-cart3},{1,-cart3},{0,cart3}] ... *)
(* Replacement rules complex and error-prone *)
\end{lstlisting}

\paragraph{Field strength substitution.}
\label{par:cs-field-strength}
\begin{lstlisting}[style=mathematica]
(* Index mismatch - doesn't apply reliably *)
F[-a,-b] -> CD[-a][A[-b]] - CD[-b][A[-a]]
\end{lstlisting}

\subsection{JSON Structure}
\label{sec:cs-json}

\subsubsection{Manual JSON (Current)}
\label{sec:cs-json-current}

File: \texttt{examples/data/chern\_simons\_3d.json}

\begin{lstlisting}[style=json]
{
  "field": "A_0",
  "rhs": {
    "terms": [
      { "coefficient": 1.0, "operator": "laplacian", "field": "A_0" },
      { "coefficient": 0.5, "operator": "gradient_x", "field": "A_2" },
      { "coefficient": -0.5, "operator": "gradient_y", "field": "A_1" }
    ]
  }
}
\end{lstlisting}

\paragraph{Key points.}
\label{par:cs-json-key-points}
\begin{itemize}
  \item CS term decomposed into gradient operators (not a composite ``curl'')
  \item Cross-field references: $A_0$ couples to $A_1$ and $A_2$
  \item Python handles this naturally via existing gradient operators
\end{itemize}

\subsubsection{Why Manual JSON}
\label{sec:cs-why-manual}

\texttt{ExportJSON.wl} operator detection issues with 3D:

\begin{itemize}
  \item Mixed \texttt{Derivative[n,m]} and \texttt{Derivative[n,m,p]} forms confuse parser
  \item Epsilon remnants in symbolic equations
  \item Cross-field detection works but coefficient extraction needs refinement
\end{itemize}

\subsection{Python Simulation}
\label{sec:cs-python}

\subsubsection{Grid Setup (2D Spatial)}
\label{sec:cs-grid}

\begin{lstlisting}[style=python]
grid = CartesianGrid(
    bounds=[(0, 50), (0, 50)],  # x and y spatial domain
    shape=[64, 64],
    periodic=True
)
\end{lstlisting}

\subsubsection{State Structure (6 Fields)}
\label{sec:cs-state}

\begin{lstlisting}[style=python]
# 3 vector components * (field + momentum)
state = FieldCollection([
    A_0, pi_0,  # Temporal component
    A_1, pi_1,  # x-spatial component
    A_2, pi_2   # y-spatial component
])
\end{lstlisting}

\subsubsection{Observed Physics}
\label{sec:cs-observed-physics}

\begin{itemize}
  \item \textbf{Initial}: Gaussian in $A_1$, others zero
  \item \textbf{Evolution}: Energy transfers to $A_0$ and $A_2$ via CS coupling
  \item \textbf{Effect}: Topological mass gives helical propagation patterns
  \item \textbf{Verification}: $\max|A_0|$ and $\max|A_2|$ grow from 0 as simulation runs
\end{itemize}

\subsection{Epsilon Tensor Component Values}
\label{sec:cs-epsilon}

For Minkowski signature $(-,+,+)$ in 2+1D:

\paragraph{Lower indices (covariant).}
\label{par:cs-lower-epsilon}
\begin{align}
  \varepsilon_{012} &= -1 \quad\text{(includes $\sqrt{|\det(g)|} = 1$ factor)} \notag \\
  \varepsilon_{021} &= +1 \notag \\
  \varepsilon_{120} &= -1 \notag \\
  &\;\;\vdots \quad\text{(all permutations with sign)} \notag
\end{align}

\paragraph{Upper indices (contravariant).}
\label{par:cs-upper-epsilon}
\begin{align}
  \varepsilon^{012} &= +1 \quad\text{(raised with metric $\eta^{\mu\nu}$)} \notag \\
  \varepsilon^{021} &= -1 \notag \\
  &\;\;\vdots \notag
\end{align}

\textbf{Critical:} Sign convention depends on metric signature. The built-in \texttt{epsiloneta3} from xAct uses standard conventions.

\subsection{Future Automation Plan}
\label{sec:cs-future}

To fully automate CS term derivation:

\subsubsection{Define Epsilon Component Rules}
\label{sec:cs-future-epsilon}

\begin{lstlisting}[style=mathematica]
epsilonRules = {
  epsiloneta3[{0,-cart3},{1,-cart3},{2,-cart3}] -> -1,
  (* ... all 6 permutations ... *)
};
\end{lstlisting}

\subsubsection{Apply After ToBasis}
\label{sec:cs-future-tobasis}

\begin{lstlisting}[style=mathematica]
componentEq = ToBasis[chart][eom]
componentEq = componentEq /. epsilonRules
componentEq = Expand[componentEq]
\end{lstlisting}

\subsubsection{Enhance ExportJSON.wl}
\label{sec:cs-future-export}

\begin{itemize}
  \item Detect curl patterns: $\ddx A_y - \ddy A_x$
  \item Handle 3-arg Derivatives consistently
  \item Map cross-derivative terms to gradient operators
\end{itemize}

\subsubsection{Test with Full Pipeline}
\label{sec:cs-future-test}

\begin{itemize}
  \item Build CS Lagrangian with epsilon
  \item \texttt{VarD} derives full EOM
  \item Decompose and export to JSON
  \item Verify coefficients and cross-field refs
\end{itemize}

\subsection{Related Examples}
\label{sec:cs-related}

\begin{itemize}
  \item \textbf{1+1D EM}: Maxwell without CS term (vector field baseline)
  \item \textbf{Coupled scalars}: Cross-field coupling pattern (scalar fields)
  \item \textbf{Chern-Simons}: Combines both: vector fields + cross-coupling + 2+1D
\end{itemize}

\subsection{References}
\label{sec:cs-references}

\subsubsection{Project Files}
\label{sec:cs-project-files}

\begin{itemize}
  \item TOML: \texttt{examples/chern\_simons/theory.toml} (derive via \texttt{tidal derive})
  \item JSON: \texttt{examples/data/chern\_simons\_3d.json}
  \item Run: \texttt{cd examples/chern\_simons \&\& bash run.sh}
\end{itemize}

\subsubsection{Physics \& Software}
\label{sec:cs-physics-refs}

\begin{description}
  \item[Deser, Jackiw \& Templeton (1982)] ``Topologically massive gauge theories.'' \emph{Annals of Physics}\ 140, 372--411. Original formulation of topologically massive gauge theory in 2+1D.
  \item[Dunne (1999)] ``Aspects of Chern-Simons Theory.'' arXiv: \href{https://arxiv.org/abs/hep-th/9902115}{hep-th/9902115}. Pedagogical review of Chern-Simons field theory.
  \item[Mart\'in-Garc\'ia et al.~\cite{martingarcia2007xact}] ``xAct: Efficient tensor computer algebra for Mathematica.'' \url{http://www.xact.es}. Symbolic tensor algebra used for the derivation pipeline.
  \item[Zwicker (2020)] ``py-pde: A Python package for solving partial differential equations.'' \emph{JOSS}\ 5(48), 2158. Original PDE backend; FD stencil conventions retained in TIDAL's native operators.
\end{description}

See \cref{sec:torsion-references} and the project's \texttt{references.bib} for the full citation list.
